\documentclass[10pt,a4paper]{article}
\usepackage{geometry}
\usepackage{geometry}
\geometry{
	a4paper,
	top=0.5cm,
	margin=2.0cm,
	headheight=1.5cm,
	includeheadfoot
 }
 
\usepackage{amsmath}
\usepackage{amsfonts}
\usepackage{amssymb}
\usepackage{graphicx}
\usepackage{lmodern}
\usepackage{enumitem}
\usepackage{gensymb}

%\usepackage[latin1]{inputenc}
\usepackage[USenglish,british,american,australian,english]{babel}

\usepackage{fancyhdr}
\usepackage{fancybox}
\usepackage{graphicx}
\usepackage[usenames, dvipsnames]{xcolor}

\usepackage{ulem}
\usepackage{lipsum}% just to generate text for the example
\usepackage{ifthen}

\usepackage{tikz}
\usetikzlibrary{
	arrows,
	automata,
        shadings,
        shadows,
        shapes,
}
\usepackage[siunitx]{circuitikz}

\usepackage[dvipsnames]{xcolor}

%%  -------------------------------------------------------------------
%%      GDS II layer, regarding MOSIS SCMOS layer map
%%  -------------------------------------------------------------------
% GDS II #41 - P_WELL
\definecolor{pwell}{rgb}{1.0, 0.74, 0.53}   % macaroni and cheese
% GDS II #42 - N_WELL
\definecolor{nwell}{rgb}{0.61, 0.87, 1.0}  % columbia blue
\definecolor{pbase}{rgb}{1.0, 0.51, 0.26}  % mango tango
\definecolor{nbase}{rgb}{0.0, 0.75, 1.0}   % capri 
% GDS II #43 - ACITVE
\definecolor{active}{rgb}{0.9, 0.4, 0.38}   % light carmine pink
% GDS II #45 - N_PLUS_SELECT
\definecolor{nimplant}{rgb}{0.45, 0.76, 0.983}% maya blue
% GDS II #44 - P_PLUS_SELECT
\definecolor{pimplant}{rgb}{1.0, 0.51, 0.26}% mango tango
% GDS II #46 - POLY
\definecolor{poly}{rgb}{0.56, 0.93, 0.56}   % light green
% GDS II #25 - CONTACT
\definecolor{contact}{rgb}{0.83, 0.83, 0.83}% light gray
% GDS II #49 - METAL1
\definecolor{metal1}{rgb}{0.38, 0.31, 0.86} % majorelle blue
% GDS II #50 - VIA1
\definecolor{via1}{rgb}{0.83, 0.83, 0.83}   % light gray
% GDS II #51 - METAL2
\definecolor{metal2}{rgb}{0.04, 0.85, 0.32} % malachite
% GDS II #61 - VIA2
\definecolor{via2}{rgb}{0.83, 0.83, 0.83}   % light gray
% GDS II #63 - METAL3
\definecolor{metal3}{rgb}{0.98, 0.93, 0.37} % maize
% GDS II #30 - VIA3
\definecolor{via3}{rgb}{0.83, 0.83, 0.83}   % light gray
% GDS II #31 - METAL4
\definecolor{metal4}{rgb}{0.75, 0.25, 0.0}  % mahogany
% GDS II #32 - VIA4
\definecolor{via4}{rgb}{0.83, 0.83, 0.83}   % light gray
% GDS II #33 - METAL5
\definecolor{metal5}{rgb}{0.79, 0.08, 0.48} % magenta (dye)
% GDS II #36 - VIA5
\definecolor{via5}{rgb}{0.83, 0.83, 0.83}   % light gray
% GDS II #37 - METAL6
\definecolor{metal6}{rgb}{0.11, 0.35, 0.02} % lincoln green
% GDS II #29 - SILICIDE_BLOCK
\definecolor{silicide-block}{rgb}{0.98, 0.94, 0.9}  % linen
% GDS II #52 - GLASS
\definecolor{glass}{rgb}{1.0, 1.0, 0.88}    % light yellow
% GDS II #26 - PADS
\definecolor{pads}{rgb}{0.75, 1.0, 0.0}     % lime (color wheel)

\definecolor{resist}{rgb}{0.71, 0.4, 0.11}  % light brown

\definecolor{silicide}{rgb}{0.29, 0.33, 0.13}
\definecolor{titanium}{rgb}{0.8, 0.58, 0.46}

\def\OpacityLayout {0.5}

%
% physical
%
\definecolor{substrate}{rgb}{0.96, 0.94, 0.93}  % isabelline
\definecolor{nitride}{rgb}{1.0, 0.03, 0.0}
\definecolor{gateoxide}{rgb}{0.88, 1.0, 1.0}    % light cyan
\definecolor{isolationoxide}{rgb}{0.84, 0.79, 0.87}% languid lavender

\usepackage{float}
\restylefloat{figure}

\setlength{\parindent}{0pt} % get rid of annoying indents

\author{David Lanzendörfer}
\title{Process datasheet: LibreSilicon 1um (hightech)}

\fancyhead[L]{
	\ifthenelse{\value{page}=1}{
		\large{\textbf{Libre Silicon 1.0$\mu m$ process:}} \\
		\textcolor{Blue}{\huge{\textbf{LS1UH}}} \\
		\small{\textbf{Hightech variant}}	 \\	
	}{
		\textcolor{Blue}{\huge{\textbf{LS1UH}}} \\
	}
}

\rhead{
	\includegraphics[width=2.0cm]{img/LSA-lightblue.png} \\
}

\addtolength{\headheight}{2\baselineskip}
\addtolength{\headheight}{0.61pt}

\pagestyle{fancy}

\renewcommand{\section}[1]{
	\vspace{20pt}
	\fcolorbox{RoyalBlue}{RoyalBlue}{
		\begin{minipage}{\textwidth}
			\textcolor{white}{\textbf{\uppercase{#1}}}
		\end{minipage}
	}
}

\renewcommand{\subsection}[1]{
	\vspace{5pt}
	\textcolor{RoyalBlue}{\textbf{\large{#1}}}
}


%\renewcommand{\headrulewidth}{0pt}

\usepackage[scaled]{helvet}
\renewcommand\familydefault{\sfdefault}
\usepackage[T1]{fontenc}

\begin{document}
\section{description}

\begin{minipage}[t]{0.45\textwidth}
	The LibreSilicons LS1UH process is a contemporary engineered 1.0 micron modular mixed signal CMOS technology for high reliability.
	Based upon the common single poly with up to 3 metal layers 1.0 micron drawn gate length and a twin-well process,
	enriched with newer features as non-volatile memory. The platform is ideal for mixed-signal applications in harsh environments,
	as well as emdedded systems
	\end{minipage}\hspace{20pt}\begin{minipage}[t]{0.45\textwidth}
	reliable over large temperature and voltage supply ranges. Fitting for many applications in infrastructure,
	industrial and consumer market. Delivered NDA-free with industrial standard design rules, precise SPICE models, analog and digital libraries.
	IPs and development kits support the process for free and open source tools.
\end{minipage}




\section{key features overview}

\begin{minipage}[t]{0.45\textwidth}
	\begin{itemize}[itemsep=0pt]
		\item 1.0-micron single poly, up to three-metal
		\item Twin-well CMOS process
		\item Lateral PNP/NPN BJT devices
		\item Integrated digital, analog and NVM in a single process
		\item Wide supply voltage range from 3.0V to 18V (max. 40V)
		\item Extended temperature range from -55$\degree C$ up to +125$\degree C$
		\item Shallow trench isolation for all MOS devices
	\end{itemize}
\end{minipage}\hspace{20pt}\begin{minipage}[t]{0.45\textwidth}
	\begin{itemize}[itemsep=0pt]
		\item Insensitive towards cosmic radiation
		\item ESD protected
		\item Low-noise PMOS/NMOS transistors
		\item High-reliability NVM using SONOS technology
		\item 7 corner case models (typical/worst/best)
		\item Community based
		\item Under free license, no NDA
	\end{itemize}
\end{minipage}

\section{applications}

\begin{minipage}[t]{0.45\textwidth}
	\begin{itemize}[itemsep=0pt]
		\item Analog and mixed-signal Circuits
		\item Avionics
		\item Infrastructure
	\end{itemize}
\end{minipage}\hspace{20pt}\begin{minipage}[t]{0.45\textwidth}
	\begin{itemize}[itemsep=0pt]
		\item Automotive
		\item Robotics
		\item Safety
	\end{itemize}
\end{minipage}

\section{primitive devices}

\begin{minipage}[t]{0.45\textwidth}
	\begin{itemize}[itemsep=0pt]
		\item PMOS/NMOS transistors
		\item SONOS flash cells
		\item Lateral PNP/NPN BJTs
	\end{itemize}
\end{minipage}\hspace{20pt}\begin{minipage}[t]{0.45\textwidth}
	\begin{itemize}[itemsep=0pt]
		\item SRAM cells
		\item Poly, Metal, Diffusion Resistors
		\item Protection, Polysilicon diodes
	\end{itemize}
\end{minipage}

\newpage

\section{ls1uh devices cross section}

\subsection{CMOS}

Below you can see the cross section of the low power MOSFETs for the CMOS logic.

\begin{figure}[H]
	\centering
	\begin{tikzpicture}[node distance = 3cm, auto, thick, scale=0.7, every node/.style={transform shape}]
		% substrate
\fill[isolationoxide] (0.0,0) rectangle (18.0,\UpperGlass);
\fill[substrate] (0,0) rectangle (18,1.25);
\node at (2,0.5) {Silicon substrate};

\fill[metal1] (1.75,2.0) rectangle (2.25,\LowerMetal);
\fill[metal1] (4.75,2.0) rectangle (5.25,\LowerMetal);
\fill[metal1] (6.0,2.0) rectangle (6.5,\LowerMetal);
\fill[metal1] (7.25,2.0) rectangle (7.75,\LowerMetal);

\fill[metal1] (10.25,2.0) rectangle (10.75,\LowerMetal);
\fill[metal1] (11.5,2.0) rectangle (12.0,\LowerMetal);
\fill[metal1] (12.75,2.0) rectangle (13.25,\LowerMetal);
\fill[metal1] (15.75,2.0) rectangle (16.25,\LowerMetal);

\fill[metal1] (1.5,\LowerMetal) rectangle (5.5,\UpperMetal);
\fill[metal1] (5.75,\LowerMetal) rectangle (6.75,\UpperMetal);
\fill[metal1] (7.0,\LowerMetal) rectangle (11.0,\UpperMetal);
\fill[metal1] (11.25,\LowerMetal) rectangle (12.25,\UpperMetal);
\fill[metal1] (12.5,\LowerMetal) rectangle (16.5,\UpperMetal);

\fill[metal2] (3.00,\UpperMetal) rectangle (4.00,\LowerMoreMetal);
\fill[metal2] (5.75,\UpperMetal) rectangle (6.75,\LowerMoreMetal);
\fill[metal2] (8.50,\UpperMetal) rectangle (9.50,\LowerMoreMetal);
\fill[metal2] (11.25,\UpperMetal) rectangle (12.25,\LowerMoreMetal);
\fill[metal2] (14.00,\UpperMetal) rectangle (15.00,\LowerMoreMetal);

\fill[metal2] (2.75,\LowerMoreMetal) rectangle (4.25,\UpperMoreMetal);
\fill[metal2] (5.50,\LowerMoreMetal) rectangle (7.0,\UpperMoreMetal);
\fill[metal2] (8.25,\LowerMoreMetal) rectangle (9.75,\UpperMoreMetal);
\fill[metal2] (11.00,\LowerMoreMetal) rectangle (12.5,\UpperMoreMetal);
\fill[metal2] (13.75,\LowerMoreMetal) rectangle (15.25,\UpperMoreMetal);

\fill[metal3] (2.75,\UpperMoreMetal) rectangle (4.25,\LowerMoreMetalTwo);
\fill[metal3] (5.50,\UpperMoreMetal) rectangle (7.0,\LowerMoreMetalTwo);
\fill[metal3] (8.25,\UpperMoreMetal) rectangle (9.75,\LowerMoreMetalTwo);
\fill[metal3] (11.00,\UpperMoreMetal) rectangle (12.5,\LowerMoreMetalTwo);
\fill[metal3] (13.75,\UpperMoreMetal) rectangle (15.25,\LowerMoreMetalTwo);

\fill[metal3] ( 2.50,\LowerMoreMetalTwo) rectangle ( 4.50,\UpperMoreMetalTwo);
\fill[metal3] ( 5.25,\LowerMoreMetalTwo) rectangle ( 7.25,\UpperMoreMetalTwo);
\fill[metal3] ( 8.00,\LowerMoreMetalTwo) rectangle (10.00,\UpperMoreMetalTwo);
\fill[metal3] (10.75,\LowerMoreMetalTwo) rectangle (12.75,\UpperMoreMetalTwo);
\fill[metal3] (13.50,\LowerMoreMetalTwo) rectangle (15.50,\UpperMoreMetalTwo);

\fill[white] ( 2.65,\UpperMoreMetalTwo) rectangle ( 4.35,\UpperGlass);
\fill[white] ( 5.35,\UpperMoreMetalTwo) rectangle ( 7.15,\UpperGlass);
\fill[white] ( 8.10,\UpperMoreMetalTwo) rectangle ( 9.90,\UpperGlass);
\fill[white] (10.85,\UpperMoreMetalTwo) rectangle (12.65,\UpperGlass);
\fill[white] (13.60,\UpperMoreMetalTwo) rectangle (15.40,\UpperGlass);

% now comes the active stuff:

% normal wells
\shade[upper left = nwell, upper right = nwell, lower right = substrate, lower left = substrate,] (1.25,\welldepthAsti) rectangle (8.25,2.0);
\shade[upper left = pwell, upper right = pwell, lower right = substrate, lower left = substrate,] (9.75,\welldepthAsti) rectangle (16.75,2.0);

\fill[gateoxide] (5.75,2) rectangle (6.75,2.2);
\fill[poly] (5.75,2.2) rectangle (6.75,3);

\fill[gateoxide] (11.25,2) rectangle (12.25,2.2);
\fill[poly] (11.25,2.2) rectangle (12.25,3);

\shade[upper left = nimplant, upper right = nimplant, lower right = nwell, lower left = nwell,] (1.5,1.8) rectangle (2.5,2);
\shade[upper left = nimplant, upper right = nimplant, lower right = pwell, lower left = pwell,] (10.0,1.8) rectangle (11.25,2);
\shade[upper left = nimplant, upper right = nimplant, lower right = poly, lower left = poly,] (11.25,2.7) rectangle (12.25,3.0);
\shade[upper left = nimplant, upper right = nimplant, lower right = pwell, lower left = pwell,] (12.25,1.8) rectangle (13.5,2);

\shade[upper left = pimplant, upper right = pimplant, lower right = nwell, lower left = nwell,] (4.5,1.8) rectangle (5.75,2);
\shade[upper left = pimplant, upper right = pimplant, lower right = poly, lower left = poly,] (5.75,2.7) rectangle (6.75,3.0);
\shade[upper left = pimplant, upper right = pimplant, lower right = nwell, lower left = nwell,] (6.75,1.8) rectangle (8.0,2);
\shade[upper left = pimplant, upper right = pimplant, lower right = pwell, lower left = pwell,] (15.5,1.8) rectangle (16.5,2);

\fill[silicide] (1.5,1.9) rectangle (2.5,2);
\fill[silicide] (4.5,1.9) rectangle (5.5,2);
\fill[silicide] (5.75,2.9) rectangle (6.75,3.0);
\fill[silicide] (7,1.9) rectangle (8,2);

\fill[silicide] (10.0,1.9) rectangle (11.0,2);
\fill[silicide] (11.25,2.9) rectangle (12.25,3.0);
\fill[silicide] (12.5,1.9) rectangle (13.5,2);
\fill[silicide] (15.5,1.9) rectangle (16.5,2);

\node at (5,-0.5) {\textbf{\large{PMOS}}};
\node at (13,-0.5) {\textbf{\large{NMOS}}};


	\end{tikzpicture}\\
	\caption{CMOS MOSFETs}
\end{figure}

The CMOS n-channel enhancement MOSFET and p-channel enhancement MOSFET transistors are being isolated from each
other with shallow trench isolation (STI) which is being filled with a low permittivity,
low temperature, low density oxide (LTO is 5 times less dense than thermal oxide).
The gate dielectricum is a 40nm thermal silicon dioxide layer which results in a break through voltage of around 40V and above.
The threshold voltage of the transistors have been dimensioned to be around +0.7V and -0.7V respectively.

\end{document}
